\vkthmsync{11}
\paragraph{\trSphericalObject{}。}\label{sub:spherical-objects}
设 $\scr{A}$ 是一个 \trCfinite{}且\trTriangulated{}的 \trACategory。沿着一个
原本出自\trKontsevich{}的想法，我们称 $Y_0\in\Ob\scr{A}$ 是一个维数为
$n\in\ZZ$ 的\textbf{\trSphericalObject}（spherical object），若它满足下面两个
条件：
\begin{itemize}
\item （\trPoincareDuality）存在一个\trLinear{}映射
$\int\colon\Hom^n_{H(\scr{A})}(Y_0,Y_0)\to\KK$，使得对一切 $Y_1$，如下%
\trPairing
\begin{equation*}
\medmuskip=1.5mu\relax\thickmuskip=3mu\relax
\Hom^{n-\bullet}_{H(\scr{A})}(Y_1,Y_0)\otimes\Hom^{\bullet}_{H(\scr{A})}(Y_0,Y_1)
\xrightarrow{\ \text{\trComposition}\ }
\Hom^{n}_{H(\scr{A})}(Y_0,Y_0)
\xrightarrow{\ \int\ }\KK
\end{equation*}
是\trNondegenerate{}的；
\item （球性\trEndomorphism{}环）$Y_0$ 的\trEndomorphism{}空间可以纳入一个%
\trShortExactSequence
\begin{equation*}
0\longrightarrow\KK\xrightarrow{\ [e_{Y_0}]\ }\Hom_{H(\scr{A})}(Y_0,Y_0)
\xrightarrow{\ \int\ }\KK\longrightarrow 0.
\end{equation*}
\end{itemize}
若 $n\ne 0$，则最后一个条件无异于说
$\Hom_{H(\scr{A})}(Y_0,Y_0)\cong\KK[t]/t^2$，其中 $t$ 的\trDegree{}为 $n$。
当 $n>0$ 时，这恰好就是 $S^n$ 的通常\trCohomology{}环，这一术语的由来即在于此。
我们日后会看到，\trLagrange{}\trSphere（以及\trHomology{}\trSphere）都会成为%
\trFukayaCategory{}中的\trSphericalObject。

\begin{lemma}\label{thm:spherical-twist-quasi-equivalence}
若 $Y_0$ 是球性的，则 $T_{Y_0}$ 是从 $\scr{A}$ 到其自身的\trQuasiequivalence。
\end{lemma}

\begin{proof}
取 $\scr{A}$ 的两个\trObject{} $Y_k$（$k=1,2$），令 $C$ 是映射
$\scr{T}^1_{Y_0}\colon\hom_{\scr{Q}}(\scr{Y}_1,\scr{Y}_2)
\to\hom_{\scr{Q}}\bigl(\scr{T}_{Y_0}\scr{Y}_1,T_{Y_0}\scr{Y}_2\bigr)$ 的%
\trMappingCone（就通常\trChainComplex{}的意义而言）。它带有一个三步的递减%
\trFiltration{} $F^\bullet C$，其相伴\trGraded{}空间为
\begin{equation*}
\medmuskip=1.5mu\relax\thickmuskip=3mu\relax
\begin{aligned}
C/F^1C&=\hom_{\scr{Q}}\bigl(\mathscr{Y}_1,
\hom_{\scr{A}}(Y_0,Y_2)\otimes\mathscr{Y}_0\bigr)[1],\\
F^1C/F^2C&=\hom_{\scr{Q}}(\mathscr{Y}_1,\mathscr{Y}_2)[1]
\oplus\hom_{\scr{Q}}(\mathscr{Y}_1,\mathscr{Y}_2)\\
&\quad\oplus\hom_{\scr{Q}}\bigl(\hom_{\scr{A}}(Y_0,Y_1)\otimes\mathscr{Y}_0,
\hom_{\scr{A}}(Y_0,Y_2)\otimes\mathscr{Y}_0\bigr),\\
F^2C&=\hom_{\scr{Q}}\bigl(\hom_{\scr{A}}(Y_0,Y_1)\otimes\mathscr{Y}_0,
\mathscr{Y}_2\bigr)[-1].
\end{aligned}
\end{equation*}
由于在 $F^1C/F^2C$ 层级上的相消，所得\trSpectralSequence{}有如下非零项：
\begin{equation*}
\medmuskip=1.5mu\relax\thickmuskip=3mu\relax
E_1^{rs}=\left\{
\begin{array}{ll}
\bigl(\Hom_{H(\scr{A})}(Y_0,Y_2)\otimes\Hom_{H(\scr{A})}(Y_1,Y_0)\bigr)^{r+s+1}
& \text{若 } r=0,\\
\bigl(\Hom_{H(\scr{A})}(Y_0,Y_2)\otimes\\
\quad\otimes\Hom_{H(\scr{A})}(Y_0,Y_0)\otimes
\Hom_{H(\scr{A})}(Y_0,Y_1)^\vee\bigr)^{r+s}
& \text{若 } r=1,\\
\bigl(\Hom_{H(\scr{A})}(Y_0,Y_2)\otimes\Hom_{H(\scr{A})}(Y_0,Y_1)^\vee
\bigr)^{r+s-1}
& \text{若 } r=2.
\end{array}\right.
\end{equation*}
$E_1$ 层级上的\trDifferential{} $\partial^{rs}$ 由 $[\phi_k]$ 诱导，这就是说，它由
$H(\scr{A})$ 中的\trComposition{}映射或其对偶组成。例如，$\partial^{1\bullet}$
是\trSurjective，因为
\begin{equation*}
\begin{tikzcd}
\Hom_{H(\scr{A})}(Y_0,Y_2)\otimes\Hom_{H(\scr{A})}(Y_0,Y_1)^\vee
\ar[d,"{\mathrm{Id}\otimes[e_{Y_0}]\otimes\mathrm{Id}}"]\\
\Hom_{H(\scr{A})}(Y_0,Y_2)\otimes\Hom_{H(\scr{A})}(Y_0,Y_0)\otimes
\Hom_{H(\scr{A})}(Y_0,Y_1)^\vee
\ar[d,"\partial^{1\bullet}"]\\
\Hom_{H(\scr{A})}(Y_0,Y_2)\otimes\Hom_{H(\scr{A})}(Y_0,Y_1)^\vee
\end{tikzcd}
\end{equation*}
是恒等映射。类似地，$\partial^{0\bullet}$ 是\trInjective，因为与%
\trPoincare{}映射\trComposition{}之后便得到一个\trIsomorphism：
\begin{equation*}
\begin{tikzcd}
\Hom_{H(\scr{A})}(Y_0,Y_2)\otimes\Hom_{H(\scr{A})}(Y_1,Y_0)
\ar[d,"\partial^{0\bullet}"]\\
\Hom_{H(\scr{A})}(Y_0,Y_2)\otimes\Hom_{H(\scr{A})}(Y_0,Y_0)\otimes
\Hom_{H(\scr{A})}(Y_0,Y_1)^\vee
\ar[d,"\mathrm{Id}\otimes{\textstyle\int}\otimes\mathrm{Id}"]\\
\Hom_{H(\scr{A})}(Y_0,Y_2)\otimes\Hom_{H(\scr{A})}(Y_0,Y_1)^\vee.
\end{tikzcd}
\end{equation*}
数维数的计算表明，该\trSpectralSequence{}的 $E_2$ 页为零。按构造，这就蕴含
$T_{Y_0}$ 是\trCohomology\trFullyFaithful{}的。

剩下要证的是它在\trObject{}的\trIsomorphism{}类上的满性。把由与某个
$T_{Y_0}(Y_1)$ \trIsomorphic{}的\trObject{}所组成的\trFull{} \trAinfty-子%
\trCategory{}记作 $\mathscr{B}\subset\scr{A}$。由
\eqref{eq:twisted-yoneda-triangle} 与 $Y_0$ 的球性可知
\begin{equation*}
T_{Y_0}\bigl(S^{n-1}Y_0\bigr)\cong Y_0.
\end{equation*}
故 $Y_0$ 落在 $\mathscr{B}$ 中。利用这一点并再次应用
\eqref{eq:twisted-yoneda-triangle}，可以看出 $\mathscr{B}$ 生成 $\scr{A}$。
另一方面，\cref{thm:afunctor-maps-exact-triangles} 表明 $\mathscr{B}$ 自身就是
一个\trTriangulated{} \trACategory，从而\trInclusionMap{}%
$\mathscr{B}\hookrightarrow\scr{A}$ 必是一个\trQuasiequivalence。
\end{proof}

\begin{remark}\label{thm:adjoint-twist-inverse}
事实上，伴随扭转 $T^\vee_{Y_0}$ 是 $T_{Y_0}$ 的一个逆，至多相差
$H^0(\mathit{fun}(\scr{A},\scr{A}))$ 中的一个\trIsomorphism。我们不打算证明
这一点（\trDerivedCategory{}语境下的平行定理见 [134, Proposition 2.10]），但我们的确
想解释一个部分性的结果，即
\begin{equation}\label{eq:adjoint-twist-composition}
T_{Y_0}T^\vee_{Y_0}(Y_1)\cong Y_1
\quad\text{在 }H^0(\scr{A})\text{ 中。}
\end{equation}
为看清这一点，过渡到\trModule{}\trCategory{} $\scr{Q}=\mathit{mod}(\scr{A})$，
并取\eqref{eq:adjoint-twist-module} 中的 $\scr{C}^\vee$。$\scr{C}^\vee(Y_0)$ 的%
\trCohomology{}处于如下的\trLongExactSequence
\begin{equation}\label{eq:adjoint-les}
\begin{tikzcd}
 & \vdots & \\
 & \Hom_{H(\scr{A})}(Y_1,Y_0)[-1] \ar[d] & \\
 & \Hom_{H(\scr{A})}(Y_0,Y_1)^\vee\otimes\Hom_{H(\scr{A})}(Y_0,Y_0)[-1]
 \ar[d] & \\
 & H(\mathscr{C}^\vee(Y_0)) & \\
 & \vdots &
\end{tikzcd}
\end{equation}
之中，其中第一个\trHomomorphism{}又是一个对偶\trComposition{}映射。据%
\trPoincareDuality，它在经过\trProjection{}
$\int\colon\Hom_{H(\scr{A})}(Y_0,Y_0)\to\KK$ 之后成为一个%
\trIsomorphism。利用 $Y_0$ 的\trEndomorphism{}环的球性，我们得出：
\eqref{eq:adjoint-les} 中第二个映射限制在子空间
$\Hom_{H(\scr{A})}(Y_0,Y_1)^\vee\otimes\KK\cdot[e_{Y_0}]$ 上是一个%
\trIsomorphism。因此，在把 $\scr{T}_{Y_0}\scr{C}^\vee$ 构造成抽象%
\trEvaluationMap{}的\trCone{}时，可以把 $\scr{C}^\vee(Y_0)$ 换成%
\trQuasiisomorphic{}的子\trComplex{}%
$\hom_{\scr{A}}(Y_0,Y_1)^\vee
=\Hom_{\KK}\bigl(\hom_{\scr{A}}(Y_0,Y_1),\KK\cdot e_{Y_0}\bigr)$。这一替换的
产物是如下 \trAinfty-模\footnote{译注：下式 $\mu^d_{\mathscr{M}}$ 行的第二个
和项原书印作 $\mu^d_{\scr{A}}(c,a_{d_1},\ldots,a_1)$；疑应为
$\mu^d_{\scr{A}}(c,a_{d-1},\ldots,a_1)$（与其前后变元的
$a_{d-1},\ldots,a_1$ 平行，且 $d_1$ 在文中无定义）；此处径改。——译者注}
\begin{equation*}
\medmuskip=1.5mu\relax\thickmuskip=3mu\relax
\begin{aligned}
\mathscr{M}(X)&=\bigl(\hom_{\scr{A}}(Y_0,Y_1)^\vee
\otimes\hom_{\scr{A}}(X,Y_0)\bigr)\oplus\hom_{\scr{A}}(X,Y_1)\\
&\quad\oplus\Hom_{\KK}\bigl(\hom_{\scr{A}}(Y_0,Y_1),
\hom_{\scr{A}}(X,Y_0)\bigr)[-1],\\
\mu^1_{\mathscr{M}}(\gamma\otimes b,c,\beta)&=\bigl((-1)^{|b|+|\gamma|}
(\gamma\circ\mu^1_{\scr{A}})\otimes b+\gamma\circ\mu^1_{\scr{A}}(b),\
\mu^1_{\scr{A}}(c),\ \mu^1_{\scr{A}}\circ\beta\\
&\quad+(-1)^{|\beta|-1}\beta\circ\mu^1_{\scr{A}}+\mu^2_{\scr{A}}(c,\cdot)
+\mu^2_{\scr{A}}(e_{Y_0},b)\otimes\gamma(\cdot)\bigr),\\
\mu^d_{\mathscr{M}}\bigl((\gamma\otimes b,c,\beta),a_{d-1},\ldots,a_1\bigr)&\\
&\quad=\bigl(\gamma\otimes\mu^d_{\scr{A}}(b,a_{d-1},\ldots,a_1),\
\mu^d_{\scr{A}}\bigl(c,a_{d-1},\ldots,a_1\bigr),\\
&\quad\ \mu^d_{\scr{A}}\bigl(\beta(\cdot),a_{d-1},\ldots,a_1\bigr)\\
&\qquad+\mu^{d+1}_{\scr{A}}(\cdot,c,a_{d-1},\ldots,a_1)\\
&\qquad+\mu^{d+1}_{\scr{A}}\bigl(e_{Y_0},b,a_{d-1},\ldots,a_1\bigr)
\otimes\gamma(\cdot)\bigr).
\end{aligned}
\end{equation*}
存在一个\trModule{}\trHomomorphism{} $\scr{M}\to\scr{Y}_1$，它把
$(\gamma\otimes b,c,\beta)$ 送到 $(-1)^{|c|}c$，更高阶的项全部为零。它的%
\trKernel{}是 $\scr{M}(X)$ 中 $c=0$ 的那个子\trComplex；利用 $e_{Y_0}$ 在%
\trCohomology{}层面上是恒等这一事实，容易看出它是\trAcyclic{}的。于是，在
$H^0(\scr{Q})$ 中我们有\trIsomorphism{} $\scr{T}_{Y_0}\scr{C}^\vee\cong\scr{M}
\cong\scr{Y}_1$。用定义换算回到 $H^0(\scr{A})$ 之后，这就建立了
\eqref{eq:adjoint-twist-composition}。
\end{remark}
